\chapter{Conclusion générale}

Ce projet a permis de construire un pipeline complet de deep learning couvrant trois régimes de données complémentaires : tabulaire, image et séquentiel. Sur les données tabulaires, le MLP s’est révélé efficace dès lors que les variables sont correctement normalisées et que l’architecture reste simple. Sur les images, les CNN ont permis d’étudier concrètement le rôle de la localité, du partage des poids, du padding, du stride et du pooling. Sur les séquences, les architectures récurrentes ont montré comment la mémoire contextuelle devient essentielle dès que l’ordre des observations conditionne la prédiction.

Au-delà des performances chiffrées, la principale conclusion est méthodologique : une bonne architecture n’est pas seulement un modèle puissant, c’est un modèle dont les hypothèses implicites sont compatibles avec la structure des données. Le deep learning apparaît alors non comme une recette universelle, mais comme une famille de méthodes dont la valeur dépend de l’ajustement entre représentation, géométrie et objectif d’apprentissage.

Le projet ouvre naturellement vers plusieurs prolongements : intégration de modèles tabulaires non neuronaux pour enrichir la comparaison, recours à des CNN plus profonds avec augmentation de données, ou encore ajout d’un mécanisme d’attention et d’un Transformer pour la traduction. Dans le cadre du module, l’ensemble obtenu constitue toutefois une base cohérente, rigoureuse et pleinement exploitable.

